Aeronautical information on airport infrastructure reaches its users through official products, such as the Aeronautical Information Publication and aeronautical charts, which present the state of the data at the effective date of each edition. 
When a value changes between editions, the change may reflect a new survey, construction work or a correction. The published product, however, states neither the cause, who asserted it, nor the premises on which the new value depends. This work addresses this traceability gap, without attempting to verify the accuracy of the measurements. It proposes an event-sourced data architecture for the airside of airports, in which every change is recorded as an event appended to a log that admits neither editing nor deletion, with its author and cause, in the manner of an accounting ledger. Values that can be computed, such as lengths and declared distances, are not stored: they are derived from their premises at query time. The architecture was implemented as an open-source library and evaluated, following design science research, through six experiments using real data from the runway of São José dos Campos Airport, the first of which characterizes the gap, with pass criteria defined before execution. Published values were recomputed within the official precision, the state of the data at past instants was reconstructed together with the cause of each change, modifications made in the middle of the history were detected, and an independent verifier confirmed its integrity without relying on the library. In the case studied, the architecture is shown to make these data traceable and verifiable, within stated limits: auditability is documentary rather than physical; a full rewrite of the history is detected only by comparison with an external reference, which was not implemented; authorship is recorded but not authenticated; and the results come from a single case study.

\bigskip
\noindent\textbf{Keywords:} Aeronautical information. Airports.
Auditability. Event sourcing. Linear referencing.